\documentclass{article}

% NeurIPS style: 9 pages of main text (plus references and appendix).
% `preprint` shows author names; use \usepackage{neurips_2025} alone for an anonymous submission.
\usepackage[preprint]{neurips_2025}

\usepackage[utf8]{inputenc}
\usepackage[T1]{fontenc}
\usepackage{hyperref}
\usepackage{url}
\usepackage{booktabs}
\usepackage{amsfonts,amsmath,amssymb,amsthm}
\usepackage{nicefrac}
\usepackage{microtype}
\usepackage{xcolor}
\usepackage{graphicx}
\usepackage{subcaption}
\usepackage{multirow}
\usepackage{listings}
\usepackage{tikz}
\usetikzlibrary{arrows.meta,positioning,fit,calc,shapes.misc}
\usepackage{pgfplots}
\pgfplotsset{compat=1.18}
\usepackage[capitalize,noabbrev]{cleveref}
\usepackage{xspace}

\hypersetup{colorlinks=true,linkcolor=blue!50!black,citecolor=blue!50!black,urlcolor=blue!50!black}

\theoremstyle{plain}
\newtheorem{theorem}{Theorem}
\newtheorem{proposition}{Proposition}
\newtheorem{corollary}{Corollary}
\theoremstyle{definition}
\newtheorem{definition}{Definition}
\newtheorem{prediction}{Prediction}
\theoremstyle{remark}
\newtheorem{remark}{Remark}

\lstdefinelanguage{nerve}{
  morekeywords={nerve,name,state,actions,transitions,invariants,ensures,futures,learn,initial,when,type,range,values,params,kind,over,after,holds,fallback,examples,action,next,machine,blocks,wires,reads,emits,bindings,on,clock,uses},
  sensitive=true, morecomment=[l]{\#}, morestring=[b]", morestring=[b]'
}
\lstset{language=nerve,basicstyle=\ttfamily\scriptsize,keywordstyle=\bfseries,commentstyle=\color{gray},
  columns=fullflexible,keepspaces=true,breaklines=true,frame=single,framerule=0.3pt,xleftmargin=4pt,xrightmargin=4pt,aboveskip=4pt,belowskip=4pt}

\newcommand{\method}{\textsc{Nerve}\xspace}
\newcommand{\Sset}{\mathcal{S}}
\newcommand{\Aset}{\mathcal{A}}
\newcommand{\Eset}{\mathcal{E}}
\newcommand{\relu}{\operatorname{relu}}
\newcommand{\Ec}{E_{\mathrm{c}}}
\newcommand{\Ej}{E_{\mathrm{J}}}
\newcommand{\Ef}{E_{\mathrm{F}}}
\newcommand{\bx}{\mathbf{x}}
\newcommand{\bz}{\mathbf{z}}
\newcommand{\R}{\mathbb{R}}
\newcommand{\ind}[1]{\mathbb{1}[#1]}
\newcommand{\step}{\mathsf{step}}
\newcommand{\have}{\textcolor{blue!50!black}{\textsf{[have]}}\xspace}
\newcommand{\plan}{\textcolor{red!50!black}{\textsf{[plan]}}\xspace}

\input{numbers}

\title{Machines Without Code\\
\large Wiring Exact Networks into Interactive Systems (\method, Part II)}

\author{%
  Alex Welcing\\
  Lupine Science, Inc.\\
  \texttt{alexwelcing@gmail.com}\\
}

\begin{document}

\maketitle

\begin{abstract}
Part I asked whether a small interactive program can be shipped as a neural network held to its behaviour
specification, and showed that it can, one interface at a time. This paper changes the question from
generalisation to exactness, reports the result, and plans the step from interfaces to machines. Trained on the
whole of its spec, one network per interface is exact: of the \nPlanned{} interfaces we set out to train,
\nNewDemos{} of them rebuilt from the demo set of declarative dragging (a sliding puzzle, Tower of Hanoi with a
solver given only by example, a kanban board, Sokoban, a meeting scheduler), every finite network that has
finished (\nExactSpecs{} of \nFiniteSpecs{} as this draft is built) reproduces every transition with no rule
broken, on every seed, at \minParamsK--\maxParamsK k parameters and minutes to an hour on a local GPU;
restricting a network's answers to the states the spec's invariants allow removes rule breaks without
retraining. The recipe stops at an
enumeration wall, and we propose to go around it rather than through it: machines wired from exact blocks. We
define blocks with ports, three couplings (independent, event-wired, read-wired) with their training cost,
argue, with a proof sketch, that exactness composes under correct wiring, and specify the block architecture (heads chosen by field
type, a proposer for large state spaces, a compiled core with trained holes), the runtime (a synchronous tick,
bindings for keyboard and mouse, selection and undo, drags inside a running machine, discrete hand-offs between
continuous blocks) and an evaluation protocol for machines. Five demos are registered as experiments with
predictions and falsifiers: an odometer of digit networks, a self-running chain checked at every hand-off, a
machine you can grab mid-run, a desktop card board under keyboard and mouse, and a ball rolling across track
blocks. Everything measured is generated from the scripts and logs committed with this paper.

\end{abstract}

\section{Introduction}
\label{sec:intro}

Software is a state machine, and Part I of this work~\citep{welcing2026nerve} asked what it takes to ship a small one as a network rather than as code: a network that takes the interface's
state and a user's action and returns the next state, held to a behaviour specification the way code is held
to types and tests. The answer was a spec language with a derived oracle, a compiler from the explicit part of
a spec to exact weights, an energy model for the rest, and a protocol under which seven small finite interfaces were
learned-only exact on cells they had never seen and a continuous slider was within $1\%$ of range. It was a result about one interface at a time, and its own
limitations section said so: nothing composed, the largest interface had $404$ states, and the only inputs
were a click and a drag.

Two things have changed since. The first is the goal. An interface's network is trained on that interface, on
the whole of its spec, so whether it would have generalised from part of the spec is beside the point; the
questions are a compiler's: exact, every time, at what cost. The second is who else is compiling machines into
weights. Qwenputer and qwenGPU~\citep{qwenputer,qwengpu} compile a 32-bit RISC-V CPU and a SIMT GPU into layers
in front of a language model of two to four billion parameters, with the model's own recurrent state as the
machine's memory, so that a general machine lives inside a model that drives it through its own tool calls.
That is one road to programs in weights, and it is not ours. Our parts are four orders of magnitude smaller,
cannot talk, and are exact by a check rather than by a construction; one of our steps is one network
evaluation, where one word of a program's output from Qwenputer is tens of thousands of instruction tokens and
tens of seconds. The question our parts raise is how they become something as large as a desktop application
without becoming one large network, which the training cost forbids.

This paper is in two halves. The first (\cref{sec:standing}) is measured: of the \nPlanned{} interfaces we set out to
train, \nNewDemos{} of them rebuilt from the demo set of declarative dragging~\citep{dragology} because that
library was built to show how much a drag can mean, one network per interface trained on its whole spec is
exact on every finite interface that has finished (\nExactSpecs{} of \nFiniteSpecs{} as this draft is built),
on every seed, with no rule broken, including a Tower of Hanoi whose solver the spec shows only by
example; restricting answers to the states the invariants allow removes rule breaks on a starved network without
retraining; and the recipe's cost grows with states times cells until it stops. The second half
(\cref{sec:blocks,sec:arch,sec:runtime,sec:demos,sec:protocol}) is a plan, written as a paper so that it can be
scored. We propose that large interactive systems be \emph{machines}: blocks that own their state and expose
ports; three couplings, independent, event-wired and read-wired, each with a training cost; an argument that
exactness composes under correct wiring; a block architecture whose heads are chosen by the types of a spec's
fields and whose candidates come from a proposer when the state space is large; a runtime with a synchronous
tick, bindings for keys, modifiers and the mouse, selection, focus and undo, drags inside a running machine, and
a rule that keeps continuous motion exact across hand-offs; and five demos, from an odometer of digit networks
to a desktop card board, each registered with a prediction and the result that would falsify it. The picture
the plan is built toward is literal: a machine you can watch run through a series of separate networks, each
hand-off checked against the spec as it happens, that you can reach into and drag while it runs.

\paragraph{Contributions.}
(1)~One network per interface, trained on its whole spec, is exact: \nPlanned{} interfaces, \nNewDemos{} from
Dragology's demo set, with generated tables, scripted provenance, and a valid-candidate rule that removes rule
breaks without retraining (\cref{sec:standing}).
(2)~A calculus of blocks, ports and wiring with a cost model for three couplings and a proof sketch that
exactness composes (\cref{sec:blocks}).
(3)~A block architecture: heads by field type, a proposer for large spaces, a compiled core with trained holes,
the energy as a port type, and explicit size and latency budgets (\cref{sec:arch}).
(4)~A runtime: the synchronous tick, a declarative input model, selection, focus and undo, drags inside running
machines, discrete hand-offs between continuous blocks, and derived rendering (\cref{sec:runtime}).
(5)~Five demos registered as experiments with predictions and falsifiers, and an evaluation protocol for
machines of networks (\cref{sec:demos,sec:protocol}).

\section{Where we stand: one network per interface is exact}
\label{sec:standing}

Part I trained on $85\%$ of an interface's (state, action) cells and scored the rest, because the question was
whether a network could generalise a rule it had seen only in part. That is not the question this programme
asks. Each interface gets its own network, trained on the whole of its spec; nothing is held out, because
nothing needs to be. The questions that remain are the ones a compiler is asked: is the output exact, is it
exact every time, and what does it cost. This section answers them for \nPlanned{} interfaces, \nNewDemos{} of them rebuilt
from the demo set of declarative dragging~\citep{dragology}, and then shows where the one-network-per-interface recipe stops.

\subsection{Trained on its whole spec, one network per interface is exact}
\label{sec:round2}

\paragraph{The interfaces.} \Cref{tab:round2} lists them. \nBefore{} were there before this round: Part I's
five hand-written specs (toggle, counter, both sliders, negotiate; its three screenshot-derived specs were never
committed and are not here) and two added since (\texttt{tracks}, a block on rails; \texttt{reorder},
Dragology's permutation demo). \nNewDemos{} are Dragology demos chosen to test what Part I's specs did not: rules whose guard reads other fields
(hanoi: only a top disk moves, never onto a smaller one); a two-dimensional reorder (kanban); a sliding puzzle
(fifteen); cars that block each other (rush hour); a rule over three objects at once (sokoban: the player, the
box ahead, the cell beyond); a list whose length changes (insert and remove); a reorder with a flag per item
(todo); meetings whose ends move without crossing (scheduler); a three-way \texttt{between} with a mode that
allows only the next place clockwise; and three continuous sliders that move together and stop together
(linked sliders). Hanoi also carries a \texttt{learn} hole: \texttt{hint}, the optimal next move toward the
goal peg, is given only by example, one per position, found by search when the spec is generated; the spec
writes the rules of the puzzle and shows, but does not write, how to solve it.

Each spec header records what was scaled down to fit enumeration (fifteen is $3\times2$, not $4\times4$;
sokoban is the top of Dragology's level; rush hour is $5\times5$ with six cars) and what Dragology does that
the spec language cannot yet say: drags that chain several snaps in one gesture, a snap radius, text, and
drawing anything but points. The eight specs written by coding agents have integration tests
(\texttt{crates/nerve/tests/demo\_*.rs}) showing that the drag futures the engine derives equal the list
Dragology's demo writes by hand (\texttt{moveItem(items, i, j)} for every $j$, the four directions a Sokoban
player may step, every stop a car can reach), from several starting states and, for fifteen, rush hour,
sokoban, insert/remove and the switch, from every valid one; hanoi, kanban, reorder and tracks have unit tests
of the same kind in \texttt{drag.rs}. The spec language has no list type, so a list is one position field per
item and an invariant that no two share a place; kanban's 360 boards live among $20{,}736$ slot tuples, and
\cref{sec:valid} is about that gap.

\begin{table}[t]
\centering
\caption{\textbf{Trained on its whole spec, one network per interface is exact.} Each row is one interface and
its own network, trained on every (state, action) cell with the configuration of Part I (300 epochs, AdamW,
weight decay $0.05$, factored head on) on a local RTX A4500; continuous specs train on $8{,}000$ sampled pairs
for $1{,}000$ epochs and are scored on $1{,}000$ fresh ones. \emph{States} counts states the invariants allow;
\emph{cells} counts legal (state, action) pairs; \emph{wall s} is wall-clock training time with several runs
sharing the GPU (\texttt{docs/paper2/timings.md}); \emph{exact} is the worst seed and \emph{violations} the
worst seed's rate of answers that break a rule. Rows are generated from \texttt{experiments/round2/results.jsonl}
and fill in as the sweep completes.}
\label{tab:round2}
\small
\setlength{\tabcolsep}{4pt}
\resizebox{\linewidth}{!}{%
\begin{tabular}{llp{4.6cm}rrrrrlr}
\toprule
interface & source & what it tests & states & cells & params & wall s & seeds & exact & viol. \\
\midrule
\input{table_round2_rows}% ends with \bottomrule
\end{tabular}}
\end{table}

\paragraph{Result.} Every finite interface is exact: $100\%$ of its cells, no violation, on every seed tried
(\nExactSpecs{} of \nFiniteSpecs{}; the seeds column says how many each has, three for all but the two
largest). The scheduler, the largest at \maxCells{} cells, was exact from its first evaluation. Hanoi reproduces all 81 solver examples on all three seeds while also being exact on its 729
written cells, so a hole and a rule can share one network of $105$k parameters without the examples being
drowned by the rules. The continuous interfaces are a different story. The two with one continuous field are
within $1\%$ of range on $99.6\%$ and $100\%$ of fresh transitions with $0.1\%$ and $0$ violations, at Part I's
tolerance under a longer recipe. The third, linked sliders, is the first interface in this programme that a
network trained on its whole spec does \emph{not} reproduce: $10.4\%$ of fresh transitions within $1\%$ of
range, a mean error of $7.5\%$ of range, and $89.6\%$ of answers breaking a rule. Its three continuous fields
move together and stop together at a track end, so the step is a descent over three coupled dimensions from
eleven seeds, and its \texttt{ensures} hold the group's shape to $0.001$, a thousandth of what the network's
tolerance allows, so an answer that is nearly right still breaks the rule. We leave the row in the table as it
is; \cref{sec:handoff} and \cref{sec:limits} draw the consequence, which is that continuous blocks must be
small and their hand-offs discrete. Nothing in \cref{tab:round2} needed a larger network than Part
I's, or a new idea; it needed exhaustive coverage of the cells and the discipline below. The whole sweep is
\totalGpuMinutes{} wall-clock minutes with several runs sharing the GPU, and the largest interface
(\maxCellsSpec, \maxCells{} cells) took over an hour of it.

\paragraph{Discipline.} Part I's own review found specs, configurations and logs behind its tables
uncommitted, and this round found one more failure of the same kind: served models whose records had been
patched by hand after training. Every model in \cref{tab:round2} is now produced by one command
(\texttt{experiments/round2/run.sh}, committed with this paper), writes its own \texttt{model.json} once and
is never edited, is scored by a second command into an append-only log, and the set the lab serves is chosen
from that log by a third. A seed is a folder; a rerun is a new row that replaces the old one in the summary
rather than a second vote. We mention this because the next sections propose machines of many networks, and a provenance discipline
that is merely good practice for one network is a precondition for debugging twenty.

\subsection{Answering only from states the spec allows removes rule breaks without retraining}
\label{sec:valid}

Reorder was the first interface whose network broke rules when data was scarce. Trained on $30\%$ of its
cells it answered $61.7\%$ of the rest correctly and broke an invariant on $26.6\%$ of them: two items in one
slot. The trainer scored every one of the $3{,}125$ position tuples as a candidate successor, and only $120$
of them are orders. The spec knows which $120$; the network was never told.

The fix is to restrict the argmin to the states the invariants allow and give every other state infinite
energy. This is the candidate-set rule of our declarative-dragging note (a state outside the derived futures
costs infinite energy), in the spirit of Part I's Theorem~1, applied to the learned term's own argmin; and it
is the first wiring primitive of \cref{sec:blocks}: a block's output is typed by its invariant. On the same
trained networks, with no retraining (every number here is from \texttt{experiments/round1b/}), reorder went from $61.7\%$
to $95.3\%$ at $70\%$ held out and from $29.8\%$ to $68.7\%$ at $90\%$, with violations falling from $27\%$ and
$63\%$ to $0.1\%$ and $5\%$ (the residue is an \texttt{ensures}, not an invariant: a legal order in which the
moved item is not where it was dropped). Retraining over the $120$ orders added nothing ($95.2\%$ against
$95.3\%$): the network already ranked real orders correctly, and impossible tuples were beating them at the
argmin. Interfaces with no invalid states score identically either way, by construction. Valid candidates are
the default from here on; \texttt{nerve eval --all-states} scores the network on its own when a paper needs
that number, and we report both wherever they differ.

\subsection{Drags are searches over futures the spec derives}
\label{sec:drags}

Every interface in \cref{tab:round2} is operated by dragging in the lab that serves it. A drag is one
minimisation per pointer frame over the futures the spec allows from the current state: the pixel gap from
the pointer to where the grabbed element would land, plus the network's energy as a tie-breaker (a candidate's
energy less the minimum over that action's candidates, infinite outside the set), with
Dragology's operators (\texttt{closest} with stickiness, \texttt{between}, \texttt{vary} through the spec's
own clamps, \texttt{far} with hysteresis) and two added for the rails demo: \texttt{each}, a \texttt{closest}
whose every branch is a \texttt{vary}, and \texttt{floating}. The futures are never listed; they are derived by
asking the oracle which actions reach which values, so a drag cannot offer an illegal drop (a buried Hanoi disk
has nowhere to go; a car stops at the first car in its way). Two facts matter for what follows. The network's
term is exactly zero on every finite interface, because the network and the spec agree on every cell, and small
on the continuous ones; in a drag the network's job is ties and holes, not geometry. And a drag is already a
search that re-plans under the hand at every frame, which \cref{sec:runtime} will lean on when the thing being
dragged is part of a running machine.

\subsection{One network per interface hits an enumeration wall}
\label{sec:wall}

\begin{figure}[t]
\centering
\input{fig_cost}
\caption{\textbf{One network per interface hits an enumeration wall.} Wall-clock seconds for 300 epochs
against the number of (state, action) cells, finite interfaces of \cref{tab:round2}, with several runs sharing
the GPU. Cost grows roughly with the number of cells; the full fifteen puzzle has about $10^{13}$ states and
no cell count that fits on this axis.}
\label{fig:cost}
\end{figure}

The trainer enumerates the state space and trains on every cell, and its wall-clock cost grows roughly with
the number of cells (\cref{fig:cost}; the candidates' embeddings are encoded once per step, so the state count
enters weakly): hanoi's 729 cells train in six minutes, kanban's $17{,}280$ in over an hour, with several runs
sharing the GPU. The interfaces we scaled down to fit are the ones
Dragology is proud of at full size: the $4\times4$ fifteen puzzle has $16!/2\approx10^{13}$ states, a real
Sokoban level has millions, and anything with a list of lists has no fixed size at all. Two exits exist. One is
engineering: score only a gesture's futures rather than every state, which turns the wall into a per-action
cost and is on the plan (\cref{sec:futures}). The other is the subject of this paper: stop building one network
per machine. A million-state odometer is six ten-state digits and five carry wires, and its network should be
one ten-state network used six times.

\section{Machines from blocks}
\label{sec:blocks}

From here to \cref{sec:demos} the paper is a plan. Nothing in \cref{sec:blocks,sec:arch,sec:runtime} is built;
\cref{sec:demos} registers, for each piece, what we expect to measure and what result would make us abandon it.
We write the plan as a paper because the design decisions are the research: once the calculus below is fixed,
building it is engineering, and if the calculus is wrong no amount of training will rescue it.

\subsection{Blocks and ports}
\label{sec:ports}

A \emph{block} is a spec, its network, its owned fields and three kinds of port. It \emph{owns} its state
fields: no other block writes them. It \emph{accepts} actions, which are its input ports; some arrive from a
user, some from other blocks. It \emph{emits} events, declared in the spec as a new clause on a transition
(\texttt{emits: \{ wrap: "old.d == 9" \}}, a predicate over the old and new state that the runtime evaluates
on the step the block just took), which are its output ports. And it may \emph{read} fields of other blocks, named in the spec as
\texttt{reads:}; a read field is part of the block's local state for every purpose except ownership. A block's
\emph{cells} are the pairs (local state, action) over its own fields, its read fields, and its accepted actions.

A block is \emph{exact} when its network's step equals its oracle's step on every cell and breaks no invariant
of the block. That is \cref{tab:round2}'s result, restated per block. Its \emph{contract} is its invariant and
its \texttt{ensures}; with valid candidates (\cref{sec:valid}) the invariant half holds by construction at the
network's output, whatever the weights, and the \texttt{ensures} half is what training is held to.

\subsection{Three couplings and what each costs}
\label{sec:couplings}

How two blocks are joined decides what each one's network has to learn.

\paragraph{Independent.} The blocks share nothing. Their energies add, so the machine's argmin splits into one
argmin per block, and the machine's cost is the sum of the blocks' costs. The scheduler of \cref{tab:round2} is
three copies of a twelve-state meeting block that happen to be trained as one $1{,}728$-state network; as a
machine it is one network of twelve states, used three times.

\paragraph{Event-wired.} A step of block $A$ emits an event that the wiring delivers as an action to block $B$.
$B$'s cells do not mention $A$; $B$ learns \texttt{inc}, not \emph{why} it was incremented. The machine's state
space is the product of the blocks' spaces, but nothing ever enumerates it: a six-digit odometer has $10^6$
states and six copies of a ten-state block. Cost is again the sum. This is the coupling hardware is built from,
a register and a carry line, and it is the one we expect to carry most of the weight.

\paragraph{Read-wired.} $B$'s rule reads $A$'s fields. $B$'s cells are now its own states times the read
fields' domains times its actions, so the cost grows with the coupling: Sokoban's box reads the player's cell
and the cell beyond, which is still small; kanban's every card reads every other card, so the board does not
split at all and its one network is the right design. The rule for a designer is the rule for a hardware
designer: keep read ports few and narrow, and move coupling into events wherever a read can be replaced by a
message. We expect that most interfaces people actually use are closer to the odometer than to kanban.

\subsection{Exactness composes}
\label{sec:theorem}

\begin{definition}[Machine]
A machine is a finite set of blocks $B_1,\dots,B_n$ with specs $S_i$, a wiring $W$ that maps each emitted
event $(i,e)$ to a finite list of deliveries $(j,a_j)$ with $a_j$ an action of $S_j$ whose parameters are
computed from $e$'s payload, and a read relation giving each block the fields it reads. A \emph{tick} applies
a finite ordered list of external actions, each to its block, and after each one processes the resulting event
queue in a fixed order, each delivery stepping its target block in the machine state current at that moment,
until the queue is empty. Emitted events are computed by the runtime from the spec's \texttt{emits} clauses on
the block's old and new states; they are derived, not learned. The
\emph{composed spec} $S_M$ is this procedure with every block stepped by its oracle; the \emph{composed
network} $N_M$ is the same procedure with every block stepped by its network.
\end{definition}

\begin{theorem}[Exactness composes]
\label{thm:compose}
If every block is exact on its cells and every tick's event cascade terminates, then $N_M(s,a) = S_M(s,a)$ for
every reachable machine state $s$ and external action $a$, and no invariant of any block is broken after any
tick.
\end{theorem}

\begin{proof}[Proof sketch]
By induction on the cascade, one external action at a time. The action steps block $i$ from a local state
(its own fields and its read fields) that is identical under $N_M$ and $S_M$ before it; exactness gives the
same successor, and the same successor gives the same emitted events, because events are evaluated from the
spec on the states and not produced by the network. Each delivery then steps its target in a machine state that, by the inductive
hypothesis, is identical under both, with the same action and parameters, so again the same successor and the
same events; the queues stay equal, and termination makes the induction finite. Invariants: every block step is
a transition of that block's oracle, which preserves its invariant, and with valid candidates the network
cannot leave the invariant's set even where it is wrong.
\end{proof}

\begin{remark}
The theorem is about equality with the composed spec, not about the composed spec being right. Wiring is code
again, and a wiring bug is a bug; the composed oracle catches it only in the sense that a test suite catches
bugs, by running the machine against its own invariants and against properties that span blocks (\texttt{the
odometer's value equals the number of increments}), exhaustively when the reachable set is small and by
sampling when it is not. Liveness and other temporal properties are outside the theorem and belong to the
bounded-LTL checker that Part I's roadmap already lists.
\end{remark}

\begin{remark}
Termination is a property of the wiring graph. An acyclic event graph terminates trivially; a cycle (two
blocks that notify each other) needs a bound per tick, and we propose that the spec declares one and the
checker enforces it, as synchronous languages require causality to be resolved at compile time~\citep{berry1992esterel,benveniste2003synchronous}.
\end{remark}

\subsection{The composed spec}
\label{sec:composed}

The spec language gains a \texttt{machine} document that names block types, instances, wires and reads.
\Cref{fig:odometer-spec} writes the odometer. Instances of one block type share one network; the composed
oracle is derived from the block oracles and the wiring, like every oracle in \method, and \texttt{nerve check}
runs the machine against itself before anything is trained.

\begin{figure}[t]
\begin{minipage}[t]{0.47\textwidth}
\begin{lstlisting}
nerve: 1
name: digit
state:
  d: {type: int, range: [0, 9]}
actions:
  inc: {}
  reset: {}
transitions:
  inc:   { d: "(d + 1) % 10",
           emits: { wrap: "old.d == 9" } }
  reset: { d: "0" }
ensures:
  - { after: inc, holds: "d == (old.d + 1) % 10" }
\end{lstlisting}
\end{minipage}\hfill
\begin{minipage}[t]{0.47\textwidth}
\begin{lstlisting}
nerve: 1
machine: odometer
blocks:
  d0: digit
  d1: digit
  d2: digit
wires:
  - on: d0.wrap
    do: d1.inc
  - on: d1.wrap
    do: d2.inc
ensures:                 # spans blocks
  - { after: d0.inc,
      holds: "value == (old.value + 1) % 1000" }
derived:
  value: "d0.d + 10*d1.d + 100*d2.d"
\end{lstlisting}
\end{minipage}
\caption{\textbf{A machine is blocks plus wiring.} Left: a ten-state digit block that emits \texttt{wrap} when
it rolls over. Right: three instances and two carry wires make a thousand-state odometer whose network is the
digit's, used three times. The machine's own \texttt{ensures} is checked by the composed oracle, exhaustively
at this size.}
\label{fig:odometer-spec}
\end{figure}

\section{Block architecture}
\label{sec:arch}

A block's network is Part I's energy model with the changes round 2 made; this section says what it is, what
has to change for blocks to be small and for lists and grids to be rules rather than tables, and what a block
must cost for a machine of them to run at frame rate.

\subsection{The energy block today}
\label{sec:energyblock}

A finite block scores a candidate successor $s'$ of $(s,a)$ by $E(s,a,s')=\Ej+\lambda_F\Ef$: the JEPA
prediction error between a predictor's output for $(s,a)$ and a target encoder's embedding of $s'$, plus the
factored per-field term whose shift kernel is equivariant along ordinal fields (Part I, \S4). It is trained
with InfoNCE over the candidate set, a hinge on spec-wrong negatives, and a variance floor, over every cell of
the spec; hole examples are oversampled twenty to one against written cells. The candidate set is the states
the invariants allow (\cref{sec:valid}), both in the loss and at the argmin, and their target embeddings are
computed once per set of weights, which is what makes a step on $3{,}125$ candidates take under a millisecond
on a CPU (\texttt{docs/paper2/timings.md}). A continuous block contrasts each true successor against the batch and sixteen near-miss perturbations,
and steps by sixty Adam iterations on the energy over its continuous fields from eleven seeds per discrete
assignment, exactly as Part I's continuous inference. The networks in \cref{tab:round2} have \minParamsK k to \maxParamsK k
parameters; three-layer GELU MLPs with $d=64$ throughout.

\subsection{Heads by field type}
\label{sec:heads}

The factored head's shift kernel made ``$+1$'' the same function at every value of an ordinal field, and Part
I's ablation showed that is what let a rule generalise from few cells. Round 2 trains on every cell, so
generalisation is not the goal, but the same property decides \emph{how many cells a rule needs} and therefore
how large a block may be. Reorder's residual errors under scarce data are mostly of one kind: a legal order in
which the moved item is not where it was dropped, because the head has no notion that the shift rule is the
same function for item $a$ as for item $e$.

We propose that the type of a field chooses the symmetry of its head, and that the spec's type system is where
that choice is made. An ordinal field keeps the shift kernel. A field that is one of several positions of
interchangeable items (reorder's $p_a,\dots,p_e$; kanban's slots) is a \emph{list} in disguise, and its head
should be permutation-equivariant: a per-item scorer that sees its own position, the moved item's old and new
positions, and a set summary of the others, with parameters shared across items~\citep{zaheer2017deepsets,lee2019settransformer}.
Under that head the rule ``if $\mathit{from}<p\le\mathit{to}$ then $p-1$'' is one function of three numbers;
the cells it needs are the distinct $(\mathit{from},\mathit{to},p)$ triples, of order $n^3$, not the
$n!\cdot n^2$ cells of the table. A field that is a position on a grid (fifteen, sokoban, rush hour) wants a
translation-equivariant head that sees displacements, not coordinates~\citep{cohen2016gcnn}. Enums keep the
plain head. The test is not generalisation for its own sake but size: a block whose rule is learned as a
function can be trained on a sample of its cells, which is the other exit from the enumeration wall.

\subsection{Candidates from a proposer}
\label{sec:futures}

The argmin over every valid state is exact and cheap at $10^3$ states and impossible at $10^{13}$. In the lab a
drag already restricts candidates to the futures a gesture can reach, derived by the oracle; a shipped block
cannot ask an oracle. But the factored head is a proposer: its per-field logits name the $k$ most likely values
of each field, and the product of those is a candidate set of $k^F$ states that the energy then ranks. For a
block whose rule changes one or two fields per action, $k=3$ covers the truth essentially always and the
candidate set is tens of states whatever the block's size. The plan is to make this the inference path for any
block over about $10^4$ states, keep the exhaustive argmin as the reference it is checked against on small
blocks, and report the proposer's recall separately so that a miss can never be mistaken for a wrong ranking.

\subsection{Compiled core, trained holes}
\label{sec:compiled}

Part I compiled every explicit transition of a finite spec into exact ReLU weights and trained only the holes
and the generalisation; its guarantee (the compiled term dominates any learned term where the spec speaks) is
the cleanest kind of exactness this programme has. Round 2 trained everything, because the Rust trainer does not
yet have the compiler. Per block the division is natural: compile the explicit transitions, train the holes
(hanoi's solver), and train the continuous parts. A compiled block has no seed to vary and nothing to score; the
digit of \cref{fig:odometer-spec} should be compiled, not trained, and the odometer experiment of
\cref{sec:demos} runs it both ways.

Blocks meet at ports, and a port needs a notion of compatibility. We take Part I's roadmap item seriously:
the energy is the type. Two blocks are compatible at a read port if the reader's expected domain and invariant
for the field match the owner's, checked statically at wiring time, and if their energies agree on the owner's
valid states for that field, checked by sampling. A block that is swapped for a retrained one must pass the
same check before the machine accepts it, which is the machine-level version of Part I's rule that the weights
that ship are the weights that were scored.

\subsection{Sizes and budgets}
\label{sec:budgets}

Blocks must stay small for three independent reasons: the trainer's cost in states times cells
(\cref{fig:cost}), the frame budget of a running machine, and the reviewer's ability to check a block
exhaustively. We fix targets rather than hopes. A finite block has at most about $2\times10^3$ valid states and
$2\times10^5$ parameters and steps in under $5$\,ms on a CPU core (under $1$\,ms measured on $3{,}125$
candidates); a continuous block steps in under $50$\,ms ($44$\,ms measured for the descent today;
\texttt{docs/paper2/timings.md}). At $60$ frames per second a tick has $16$\,ms; a cascade of ten finite
blocks fits in it on one core, and blocks that are not on the same cascade path step in parallel. Anything
that does not fit these budgets is two blocks.

\section{Runtime}
\label{sec:runtime}

A machine of blocks needs what a single interface did not: time, input richer than a click or a drag, and a
loop that moves events between blocks while checking every step. This section specifies that loop and the input
model, and fixes the one design rule that keeps continuous motion from wandering.

\subsection{The tick}
\label{sec:tick}

The runtime is one loop. Each frame it collects inputs, turns them into actions through the bindings
(\cref{sec:bindings}), applies each action to its block, drains the event queue through the wiring with each
delivery stepping its target block, runs every stepped block's spec check, and renders. A tick is a pure
function of the machine state and the frame's inputs, so a session is a log of inputs and can be replayed
exactly, which is also how it is tested. Events emitted during a tick are processed within it, in a fixed order,
to the bound the machine declares; the clock is a block like any other, whose only action is \texttt{advance} and
whose event the machine wires to whatever it drives. This is the synchronous hypothesis of
Esterel and Lustre~\citep{berry1992esterel,halbwachs1991lustre}: a reaction is instantaneous and deterministic,
and causality is settled when the machine is checked, not when it runs. Harel's statecharts~\citep{harel1987statecharts}
and the Elm architecture~\citep{czaplicki2013elm} are the same discipline for interfaces, and ours differs only
in what steps a block.

\subsection{Bindings: the input model}
\label{sec:bindings}

Desktop software is mostly keys, modifiers and selection, not drags. None of that is new state; it is a mapping
from input events to actions, and it belongs in the spec rather than in code. A \texttt{bindings} section maps
a key or mouse event, with modifiers, to an action whose parameters come from the pointer, the focused item or
the selection (\cref{fig:bindings}). Modifiers change which action a gesture performs: Alt-drag performs
\texttt{copy} where drag performs \texttt{move} (Dragology's \texttt{alt-to-copy} and \texttt{drag-to-copy});
Shift-drag constrains a two-dimensional \texttt{vary} to one axis by fixing the other field; the wheel steps a
field; a double-click opens an item; Esc leaves a mode. The runtime owns the mapping and the blocks never see a
key, so the same machine runs under a different keyboard layout, or under a touch screen, by changing one
section.

\begin{figure}[t]
\begin{lstlisting}
bindings:
  Tab:           focus.next
  shift+Tab:     focus.prev
  ArrowRight:    board.move(card: focus.id, dx: 1, dy: 0)
  Delete:        board.remove(card: selection.items)
  ctrl+a:        selection.all
  Escape:        [selection.none, menu.close]
  click(card):   selection.only(card)
  ctrl+click:    selection.toggle(card)
  shift+click:   selection.range(focus.id, card)
  drag(card):    board.drag_card            # the gesture in `futures`
  alt+drag:      board.drag_copy
  shift+drag:    board.drag_card(axis: lock)
  wheel(board):  board.scroll(by: delta)
  rightclick:    menu.open(at: pointer, for: card)
  ctrl+z:        undo                       # runtime, not a block
\end{lstlisting}
\caption{\textbf{Desktop input is state plus bindings.} The input model of the desktop board of
\cref{sec:demos}: every key and mouse event names an action of some block, with parameters drawn from the
pointer, the focused item or the selection. Blocks never see a key.}
\label{fig:bindings}
\end{figure}

\subsection{Selection, focus, menus and undo}
\label{sec:selection}

Selection over $n$ items has $2^n$ states and would double the board's spec for every item; it is a block of its
own, read-wired to the board only through the item ids it names, and the board's actions take an item set as a
parameter. Focus is a block with one field. A context menu is a mode field with \texttt{open} and
\texttt{close} and a guard on every other action while it is open. Undo is not a block at all: it is a runtime
wrapper that keeps a history of machine states and restores one, which is exact by construction and costs
nothing to learn; learning it would multiply every block's state space by the depth of the history. Marquee
selection is a new kind of drag whose futures are selection sets: the candidates are the sets of items whose
rendered positions fall inside the rectangle swept so far, derived from the render block as other futures are
derived from the oracle, and it is the one piece of the desktop model that none of the Dragology demos we
catalogued contains. Text
entry stays out until the spec language has a text type.

\subsection{Drags inside a running machine}
\label{sec:livedrag}

A drag is a search over the futures the spec allows from the current state (\cref{sec:drags}). Inside a
running machine the grabbed block is detached from the clock for the duration of the gesture: events addressed
to it queue, the rest of the machine keeps running, and the pointer frame is minimised as before. On release
the drop is an external action and the queued events are delivered, so the cascade resumes from the dropped
state; the composed oracle says where the machine should be $k$ ticks later, and \cref{sec:demos} measures
whether it gets there. The visible signature of the whole design is that every block's spec check runs every
tick: a block that answered a legal transition lights green, one that did not lights red and the machine stops
there, with the transition on screen. Code shows nothing when it is right; this runtime shows its checks.

\subsection{Continuous hand-offs}
\label{sec:handoff}

Continuous blocks with one continuous field are within $1\%$ of range on $99.6$--$100\%$ of transitions, which
is not exact; one with three coupled fields (linked sliders, \cref{sec:round2}) is within tolerance on a tenth
of them; and the error compounds across steps: on the continuous slider a value that should stay at $0$ drifted to $0.59$ and then
$1.41$ over two clicks that the spec says do nothing (observed in the lab with the round-1 model;
\texttt{docs/paper2/timings.md}). A chain of continuous blocks would wander. The design rule
is that a hand-off is a discrete event: a continuous block emits an event only when a quantised condition holds
(the ball reached the end of the track), the event's payload is the spec's canonical value, and the receiving
block starts from that value, not from the sender's float. Inside a block the motion is continuous and may
carry its error for the length of one episode; across a hand-off the machine is exact again. The rolling-ball
experiment of \cref{sec:demos} measures drift with and without the rule.

\subsection{Rendering}
\label{sec:render}

The lab draws every element as a point, which was enough for a dial and a knob and is not enough for a
machine anyone would watch. The render block grows the shapes Part I's roadmap lists: rectangles, lines,
angles and z-order, each an expression over state, derived and never learned. The lab is the runtime: one page,
one tick loop, every block a network served from one process, with the Worker export of Part I as the way to
ship a machine once it runs. The networks are the behaviour; the picture is a function of the state, and we keep
it that way so that a machine can be checked without being looked at.

\section{The demos as experiments}
\label{sec:demos}

Each demo below is the acceptance test of one part of the plan. For each we fix, before building it, what we
expect to measure and what result would make us drop or redesign that part; the predictions are numbered so
that the next paper can report against them, as \texttt{experiments/round2/README.md} does for round 2. They
are ordered by what each one depends on, and by cost: the first is days, the last is open.

\subsection{An odometer of digit networks}
\label{sec:demo-odometer}

\emph{Setup.} The digit block and machine of \cref{fig:odometer-spec}, once compiled and once trained; three
digits checked exhaustively over all $1{,}000$ states and both actions; six digits sampled over $10^5$ random
increment sequences of length up to $10^4$ and compared with integer arithmetic.

\begin{prediction}
\label{pred:odometer}
Both odometers are exact: no mismatch with the composed oracle and no violation on any of the $2{,}000$ cells
of the three-digit machine or on any sampled sequence of the six-digit one, and the trained digit's network is
the same $\sim80$k-parameter network for every instance.
\end{prediction}
\emph{Falsifier.} Any carry error. A carry error would mean \cref{thm:compose}'s premise failed in practice
(a block exact on its cells is not exact when the action arrives as an event), and the wiring would need a
different semantics before anything larger is built.

\subsection{A machine that runs itself}
\label{sec:demo-chain}

\emph{Setup.} A clock block drives a chain of finite blocks: a toggle flips on every advance; its flip to
\texttt{on} increments a digit (the odometer's block); the digit's wrap moves an item one place in a reorder
list; the list returning to its initial order turns the three-way switch one place clockwise; the switch
coming back to \texttt{r} resets the digit. Three networks from \cref{tab:round2}, the digit, and one clock,
five wires, rendered as the machine it is, with every hand-off checked live. Continuous blocks join in
\cref{sec:demo-ball}, not here, so that agreement with the composed oracle is exact rather than within a
tolerance.

\begin{prediction}
\label{pred:chain}
The chain runs for $10^4$ ticks with no violation at any block and no disagreement with the composed oracle,
at under $16$\,ms per tick on a laptop CPU.
\end{prediction}
\emph{Falsifier.} A violation anywhere, or a tick that cannot fit the frame. The former says a hand-off delivered
a state a block was not trained for (a gap in the composed spec's reachable set that the block's spec did not
cover); the latter says the budgets of \cref{sec:budgets} are wrong.

\subsection{Grab anything while it runs}
\label{sec:demo-grab}

\emph{Setup.} The same chain; at random ticks a scripted pointer grabs a list item or the switch's knob, drags
it along a recorded path and releases; the machine's state $k$ ticks after release is compared with the
composed oracle's prediction from the dropped state, for $k\in\{1,10,100\}$ over $1{,}000$ grabs.

\begin{prediction}
\label{pred:grab}
Agreement is $100\%$ at every $k$: a drag mid-run is just another external action to the composed spec, and the
blocks downstream re-plan from the dropped state exactly.
\end{prediction}
\emph{Falsifier.} Any disagreement, which would localise to either the detach-and-queue semantics of
\cref{sec:livedrag} or to a drag future the runtime offered that the composed spec does not.

\subsection{The desktop board}
\label{sec:demo-board}

\emph{Setup.} Cards on a grid; a board block, a selection block, a focus block, a menu block and the undo
wrapper, under the bindings of \cref{fig:bindings}. Tests are replayed traces: $1{,}000$ random input sequences
of length $50$ drawn from the binding table (keys with modifiers, clicks, drags, marquees), replayed through the
runtime and through the composed oracle.

\begin{prediction}
\label{pred:board}
Runtime and composed oracle agree on every trace, and the machine answers every input within the frame.
Selection and focus as separate blocks keep the board's network no larger than kanban's.
\end{prediction}
\emph{Falsifier.} Any divergence (a bug in bindings, selection semantics or the menu's guards), or a board
network that only reaches exactness when selection is folded into it, which would say the read-coupling
through item sets is tighter than \cref{sec:couplings} assumes.

\subsection{A ball along track blocks}
\label{sec:demo-ball}

\emph{Setup.} Three continuous track blocks in a chain; a ball's position is continuous within a block; the
hand-off from one track to the next is an event with and without the snapping rule of \cref{sec:handoff}.
Measure the error in position against the composed oracle after $1$, $10$ and $100$ hand-offs over $200$ runs.

\begin{prediction}
\label{pred:ball}
Without snapping, error grows roughly linearly in the number of hand-offs and exceeds $1\%$ of range within
ten; with snapping it is bounded by one block's within-episode error ($\le1\%$) at every count.
\end{prediction}
\emph{Falsifier.} Growth with snapping on, which would mean the error enters inside a block's episode faster
than an episode ends, and the continuous block itself needs the compiled-core treatment of \cref{sec:compiled}
or a shorter episode.

\subsection{Order and cost}
The odometer first, because it tests the theorem with the least machinery and needs only \texttt{emits},
\texttt{machine} and the composed oracle. Then shapes in the render block and the chain, which is the first
thing a viewer can watch; then grabbing, which needs nothing new beyond the clock detach; then the board, which
needs bindings, the selection and focus blocks and the undo wrapper; then the ball, which needs more than one
continuous field per block and the snapping rule. Each is a demo in the lab on the day it passes its test.

\section{Evaluation protocol for machines}
\label{sec:protocol}

Part I's protocol scored one network on one interface: held-out cells, learned-only scoring, violation rate,
scores from the weights that ship, runtime parity. A machine needs a protocol of its own, and round 2's review
of our own practice added rules about evidence. We fix both here so the next paper is held to them.

\paragraph{Conformance.} Exhaustive when the machine's reachable state set is at most about $10^4$: every
reachable state, every external action, runtime against composed oracle. Sampled otherwise: random external
action sequences from the initial state, with the sampler biased toward guards and range ends as Part I's
adversary policy was. Both report exact match and violation rate, per block and for the machine.

\paragraph{Rollouts, not steps.} A single step hides accumulated error; the unit of measurement is a
$k$-step rollout, reported at $k\in\{1,10,100\}$, for discrete machines as agreement and for continuous ones as
error against the composed oracle after $k$ hand-offs.

\paragraph{Traces for interfaces with input.} A desktop interface is tested by replaying input traces, random
and recorded, through the runtime and the composed oracle; the trace is the test case, is committed with its
expected final state, and is replayed in CI.

\paragraph{Two numbers, always.} Where the spec's invariants mask the network's answer
(\cref{sec:valid}) we report the network on its own beside the network with the mask; they answer different
questions and one must never stand in for the other.

\paragraph{Seeds and evidence.} At least three seeds for any trained block in a table, the worst seed being the
reported one; this binds from the next paper on (\cref{tab:round2} has one to three seeds per row, as its seeds
column says, and its two largest interfaces were trained once). Every model is produced by a committed command, writes its record once, and is scored into an
append-only log by a second committed command; served sets are chosen from the log by a third; no artifact is
edited after the fact. Predictions are registered in the repository before a run, with the result that would
falsify them, and results are reported next to the predictions whether they held or not. Round 2's first
predictions did not hold: before the goal changed we predicted that a solver learned from examples would be
right on $85\%$ of unseen positions and measured $13\%$, and \texttt{experiments/round2/README.md} records both.

\section{Related work}
\label{sec:related}

\paragraph{Programs in weights.} Constructing automata in networks is classical~\citep{omlin1996constructing,siegelmann1995computational};
Tracr compiles RASP programs into transformer weights~\citep{lindner2023tracr}, Cajal compiles a typed linear
language to neurons with a correctness proof~\citep{cajal,cajalrec}, and ink's \texttt{emit-circuit} lowers
functions over finite types to exact ReLU nets~\citep{inklang}; Part I's compiler is the table-sized member of
this family. Qwenputer and qwenGPU go the other way in scale: a 32-bit RISC-V CPU, and a SIMT GPU, compiled into
layers in front of a pretrained language model of two to four billion parameters, with the model's own recurrent
state as RAM~\citep{qwenputer,qwengpu}. Those are general machines that a language model drives through its own tool calls and whose results it
reads back; a Qwenputer instruction is a few trace tokens, and one generated word of a program's output costs
tens of thousands of them and tens of seconds. Ours are specific machines built from exact parts of about $10^5$ parameters that step in
milliseconds and cannot talk. The two meet at Part I's roadmap item 20, a \method block compiled into layers in
front of a small Qwen, and nothing in this paper is in tension with it.

\paragraph{Synchronous languages, statecharts, and the Elm architecture.} The tick of \cref{sec:tick} is the
synchronous hypothesis of Esterel and Lustre~\citep{berry1992esterel,halbwachs1991lustre,benveniste2003synchronous}:
a reaction is instantaneous, deterministic, and its causality is settled before it runs. Statecharts~\citep{harel1987statecharts}
gave interfaces hierarchy and broadcast events; the Elm architecture~\citep{czaplicki2013elm} gave them a pure
update function over messages. Our machines borrow all three ideas and change one thing: the update function of
each component is a network held to a spec, so composition has to be argued (\cref{thm:compose}) rather than
assumed.

\paragraph{Declarative dragging.} Dragology~\citep{dragology} describes a drag by the set of states it may lead
to and gives operators (\texttt{closest}, \texttt{between}, \texttt{vary}, \texttt{whenFar}, floating, chaining)
whose semantics we restated as energy minimisation in Part I's note and implemented in round 2. Ten of its demos
are now \method specs; what we add is that the futures are derived from the spec and checked, and what we lack
is chaining, snap radii, text, and its drawing.

\paragraph{Neural interfaces and world models.} NeuralOS~\citep{rivard2026neuralos}, GameNGen~\citep{valevski2024gamengen}
and Genie~\citep{bruce2024genie} learn an interface or a game as pixels from pixels, and are impressive where
rules are soft; they have no spec to be held to and no way to say which pixel is a violation. We learn the
behaviour and derive the picture, and the question we ask of a network is a compiler's question.

\paragraph{Energy-based inference, JEPA, and equivariance.} Scoring candidates and taking the minimum is
energy-based learning~\citep{lecun2006tutorial}; the predictor-in-latent-space form is JEPA~\citep{lecun2022path,assran2023ijepa},
trained here contrastively~\citep{oord2018cpc} with a variance floor~\citep{bardes2022vicreg}. The heads of
\cref{sec:heads} are the standard constructions for sets and grids~\citep{zaheer2017deepsets,lee2019settransformer,cohen2016gcnn};
what is new is letting a spec's field types choose among them.

\paragraph{Modular networks.} Neural module networks assemble a network per question from reusable
modules~\citep{andreas2016nmn}; mixtures of experts route among sub-networks~\citep{shazeer2017moe}. Both compose
learned parts by learned routing. Our blocks compose by a wiring that is written, checked, and not learned, and
each part is exact on its own contract before it is wired; the composition result follows from that and would
not from routing.

\paragraph{Specifications and checking.} TLA\textsuperscript{+}~\citep{lamport2002tla} and Alloy~\citep{jackson2012alloy}
check machines against properties by exhaustive and bounded search; the composed oracle of \cref{sec:composed}
is the same instrument pointed at a network. Typed holes~\citep{omar2019hazel} and sketching~\citep{solar2008sketch}
are the ancestors of the \texttt{learn} holes that hanoi's solver fills.

\section{Limitations, and what we do not know}
\label{sec:limits}

\emph{Most of this paper is a plan.} \Cref{sec:standing} is measured; \cref{sec:blocks,sec:arch,sec:runtime} are
designs; \cref{sec:demos} is predictions. We have built none of the machines. The value of writing the plan as
a paper is that the next one can be scored against it, and we expect some predictions to fail: round 2's first
set did.

\emph{Exactness composes only with correct wiring.} \Cref{thm:compose} moves the burden from the networks to
the wiring and the composed spec, which are code and can be wrong in the ordinary way. The composed oracle
checks them against properties a person wrote, and a property nobody wrote is not checked. Temporal properties
are outside the theorem entirely.

\emph{Read coupling does not split.} A machine whose blocks all read each other is one block, and kanban shows
that such interfaces exist. The head programme of \cref{sec:heads} is a bet that most of them have a symmetry
that makes the one block small; it is a bet.

\emph{Continuous blocks are not exact, and coupled ones are not close.} $99.6$--$100\%$ within $1\%$ of range
is what one continuous field gives; three coupled fields gave $10.4\%$ (linked sliders). The snapping rule
contains the error between hand-offs and does nothing about it inside a block, and a block like linked sliders
needs either a different inference (one descent per field, or the compiled core of \cref{sec:compiled} for the
coupling) or to be split into three blocks wired by events, which is the design \cref{sec:couplings}
recommends anyway.

\emph{One GPU, one configuration.} \Cref{tab:round2} and \cref{fig:cost} are one machine and one training
recipe, with three seeds for the small interfaces and one for the two largest; the times are wall-clock with
several runs sharing the GPU.

\emph{Scale comparisons are approximate.} Qwenputer's figures are from its repository's README; qwenGPU's
repository was not reachable when we checked, and its size is taken from Part I's citation. The comparison is
of kinds, not of a benchmark.

\emph{What \method still cannot say.} No text, no drawing beyond points, no drag that chains several snaps,
no snap radius, no list type: Dragology's largest demos (its permutation lab, braids, tableaux) remain out of
reach, and the full-size fifteen and Sokoban levels wait on the proposer of \cref{sec:futures}.

\emph{One machine, one user, one clock.} Nothing here is about machines that talk to machines, about two
users in one session (the negotiation game's real setting), or about concurrency.

\section{Conclusion}
\label{sec:conclusion}

One network per interface, trained on the whole of its spec, is exact at about $10^5$ parameters and minutes
to an hour on a local GPU, on every finite interface that has finished, including ten from a library built to show how rich dragging can
be. That result closes the question Part I opened and opens the one this paper is about: what a large
interactive system made of such networks should look like. Our answer is a machine, not a model: small exact
blocks that own their state, wired by events and few reads, driven by a synchronous tick, operated through
bindings, checked at every step and shown being checked. Exactness composes under that discipline by a short
argument, and the cost of a machine is the sum of its parts. The demos of \cref{sec:demos} are how we will find
out whether the argument survives contact with a clock, a pointer and a keyboard; their predictions are
written down, and the next paper reports against them. If they hold, the way to ship a complex interface as
weights is not a larger network but a netlist of small exact ones, and the picture we are building toward, a
machine you can watch run through a series of separate networks and reach into while it runs, is a description
of the runtime, not a metaphor.


\paragraph{Reproducibility.} Every number in \cref{sec:standing} is generated from
\texttt{experiments/round2/results.jsonl} by \texttt{docs/paper2/make\_tables.py}; the training commands are the
committed scripts in \texttt{experiments/round2/}, and the served models are built from the same log by
\texttt{serve\_set.py}. Each spec in \cref{tab:round2} has tests in \texttt{crates/nerve/tests/} showing its drag
futures equal the Dragology rule it was taken from, computed independently.

\paragraph{Use of language models.} As in Part I, the branch and this manuscript were written by Claude (Anthropic)
acting as a coding agent under the author's direction, and the method itself uses coding agents for intake and
for the spec-writing step of \cref{sec:standing}. The author set the questions, the goal change recorded in
\cref{sec:standing}, and the plan in \cref{sec:blocks,sec:arch,sec:runtime,sec:demos}, and reviewed every claim.

\bibliographystyle{plainnat}
\bibliography{refs}

\end{document}
